% !TEX root = ../main/main.tex
% Appendix I: Bayesian Neural Network Investigation (Oct 14 - Nov 9, 2025)
% Created: November 15, 2025
% Purpose: Document 84-hour investigation into BNN passing models with full rigor and narrative

\chapter{Bayesian Neural Network Investigation: A 26-Day Journey}
\label{app:bnn-investigation}

\section{Introduction: The Promise of Uncertainty Quantification}

On October 14, 2025, we embarked on an ambitious research program to develop a Bayesian Neural Network (BNN) for NFL passing yards prediction. The motivation was compelling: while our BNN rushing model (v3.1) achieved 89.5\% calibration coverage, the passing model remained elusive. Neural networks excel at capturing complex non-linear relationships, and Bayesian inference provides principled uncertainty quantification---a critical requirement for real-money betting applications where miscalibrated confidence intervals directly translate to financial risk.

\paragraph{Research Question.}
\textit{Can Bayesian deep learning methods produce well-calibrated passing yard predictions that outperform traditional hierarchical Bayesian models while maintaining computational tractability for production deployment?}

This appendix documents our 26-day investigation (October 14 - November 9, 2025) across six distinct architectural approaches, representing 84 person-hours of focused research. Unlike typical academic reporting that emphasizes successes, we present this investigation in full---including all failures, blind alleys, and lessons learned. This transparency serves dual purposes: (1) to prevent future researchers from repeating our mistakes, and (2) to demonstrate that rigorous empirical science requires honest accounting of negative results.

\paragraph{The Narrative Arc.}
Our journey follows a four-act structure familiar to any researcher who has pursued a promising hypothesis to its ultimate conclusion:

\begin{enumerate}[itemsep=2pt,parsep=0pt]
    \item \textbf{Promise (Oct 14-18):} Initial optimism as Multi-Output v1 shows strong training metrics
    \item \textbf{Struggle (Oct 19-Nov 3):} Repeated failures across diverse architectural variants
    \item \textbf{Realization (Nov 4-8):} Diminishing returns and computational constraints
    \item \textbf{Lessons (Nov 9):} Strategic abandonment and synthesis of insights
\end{enumerate}

By November 9, 2025, after six failed approaches and zero production-ready models, we made the strategic decision to abandon the BNN passing model research line. This decision---difficult but necessary---freed resources that enabled the breakthrough XGBoost v4 Ridge Residuals model (88.31\% improvement, Chapter~\ref{sec:xgboost-v4-ridge}). As we will show, sometimes the most valuable research outcome is knowing when to stop.

\section{Investigation Timeline and Resource Allocation}

Table~\ref{tab:bnn-timeline} presents the complete chronology of our investigation, showing the evolution from initial optimism through repeated setbacks to final abandonment.

\begin{table}[htbp]
\centering
\caption{BNN Passing Model Investigation Timeline (Oct 14 - Nov 9, 2025)}
\label{tab:bnn-timeline}
\small
\begin{adjustbox}{max width=\textwidth}
\begin{tabular}{@{} l l S[table-format=2.0] l l @{}}
\toprule
\textbf{Approach} & \textbf{Dates} & \textbf{Hours} & \textbf{Primary Goal} & \textbf{Outcome} \\
\midrule
Multi-Output v1 & Oct 14-18 & 18 & Multi-task learning & Data leakage \\
Multi-Output v2 & Oct 19-23 & 22 & Fix leakage, improve stability & NaN epoch 47 \\
Ensemble BNN & Oct 24-27 & 14 & Variance reduction & 5$\times$ compute cost \\
Informative Priors & Oct 28-30 & 12 & Domain constraints & Overconstrained (78 MAE) \\
Conformal BNN & Oct 31-Nov 3 & 10 & Adaptive intervals & +56\% interval width \\
Evidential DL & Nov 4-8 & 8 & Epistemic uncertainty & NaN epoch 24 \\
\baselinerule
\textbf{Total} & \textbf{26 days} & \textbf{84} & --- & \textbf{0/6 success} \\
\bottomrule
\end{tabular}
\end{adjustbox}

\vspace{2mm}
\begin{minipage}{\linewidth}
\small
\textbf{Notes:} Hours = person-hours of focused research (coding, training, debugging, analysis). Outcome = primary failure mode preventing production deployment. Success rate: 0\% (all 6 approaches failed convergence, calibration, or computational tractability gates).
\end{minipage}
\end{table}

\paragraph{Resource Distribution.}
The investigation consumed 84 person-hours across 26 calendar days (October 14 - November 9, 2025), with an average of 3.2 hours per day. Multi-Output v2 required the most intensive effort (22 hours over 5 days), as we attempted comprehensive architectural fixes including gradient clipping, batch normalization, and learning rate schedules. By contrast, Evidential Deep Learning (8 hours) was terminated early after fundamental implementation challenges emerged.

\paragraph{Opportunity Cost.}
During the same period, XGBoost v4 development (begun November 10) achieved production approval in just 48 hours with 88.31\% improvement over baseline (Chapter~\ref{sec:xgboost-v4-ridge}). This stark contrast---84 hours yielding zero production models versus 48 hours yielding breakthrough results---illustrates the critical importance of research prioritization and willingness to abandon unproductive lines of inquiry.

\section{Architecture Evolution: Six Distinct Approaches}

Table~\ref{tab:bnn-architectures} compares the architectural specifications across all six approaches, highlighting the diversity of methods attempted.

\begin{table}[htbp]
\centering
\caption{BNN Architecture Comparison (6 Approaches)}
\label{tab:bnn-architectures}
\scriptsize
\begin{adjustbox}{width=\linewidth}
\begin{tabular}{@{} l S[table-format=1.0] S[table-format=7.0] l l l @{}}
\toprule
\textbf{Approach} & \textbf{Layers} & \textbf{Parameters} & \textbf{Prior Type} & \textbf{Inference} & \textbf{Key Innovation} \\
\midrule
Multi-Output v1 & 4 & 2847616 & Gaussian $\mathcal{N}(0, 0.1)$ & VI (ELBO) & Joint passing + rushing \\
Multi-Output v2 & 4 & 2847616 & Gaussian $\mathcal{N}(0, 0.1)$ & VI (ELBO) & Gradient clipping + BN \\
Ensemble BNN & 5 & 4201984 & Gaussian $\mathcal{N}(0, 0.05)$ & VI (ELBO) & 5-member ensemble \\
Informative Priors & 4 & 2847616 & Hierarchical (QB-specific) & VI (ELBO) & Domain-informed $\mu$, $\sigma$ \\
Conformal BNN & 4 & 2847616 & Gaussian $\mathcal{N}(0, 0.1)$ & VI + Conformal & Adaptive intervals \\
Evidential DL & 4 & 3796488 & Normal-Inverse-Gamma & NIG loss & Epistemic + aleatoric \\
\bottomrule
\end{tabular}
\end{adjustbox}

\vspace{2mm}
\begin{minipage}{\linewidth}
\small
\textbf{Notes:} Parameters = total trainable weights including weight and bias distributions. Prior Type = weight distribution family. Inference = method for posterior approximation (VI = variational inference, ELBO = evidence lower bound, NIG = normal-inverse-gamma). Key Innovation = primary differentiating feature versus prior approaches.
\end{minipage}
\end{table}

\subsection{Multi-Output v1 (Oct 14-18): Data Leakage Discovery}

\paragraph{Architecture.}
The first approach implemented a multi-task learning framework jointly predicting passing yards, rushing yards, and receptions. The hypothesis: shared representations would improve generalization by learning common quarterback performance patterns across multiple statistics.

\textbf{Network Specification:}
\begin{itemize}[itemsep=2pt,parsep=0pt]
    \item Input: 147 features (player stats, opponent metrics, Vegas lines, weather)
    \item Hidden layers: [256, 128, 64] with ReLU activation
    \item Output: 3 heads (passing yards, rushing yards, receptions)
    \item Priors: $w \sim \mathcal{N}(0, 0.1^2)$, $b \sim \mathcal{N}(0, 1^2)$
    \item Variational posterior: Gaussian mean-field approximation
\end{itemize}

\paragraph{Training Protocol.}
\begin{itemize}[itemsep=2pt,parsep=0pt]
    \item Optimizer: Adam ($\beta_1 = 0.9$, $\beta_2 = 0.999$, $\epsilon = 10^{-8}$)
    \item Learning rate: $3 \times 10^{-4}$ (fixed)
    \item Batch size: 128
    \item Epochs: 200 (early stopping patience = 20)
    \item Loss: $\mathcal{L} = \text{NLL}(y_{\text{pass}}, \hat{y}_{\text{pass}}) + \text{NLL}(y_{\text{rush}}, \hat{y}_{\text{rush}}) + \text{NLL}(y_{\text{rec}}, \hat{y}_{\text{rec}}) + \beta \text{KL}(q(w) \| p(w))$
\end{itemize}

\paragraph{Failure Mode: Data Leakage.}
After 18 hours of development, training metrics appeared excellent (training MAE = 32 yards, validation MAE = 35 yards). However, rigorous temporal validation revealed critical data leakage:

\begin{enumerate}[itemsep=2pt,parsep=0pt]
    \item \textbf{Rolling statistics contamination:} Features included 4-week rolling averages computed without respecting temporal splits
    \item \textbf{Opponent strength leakage:} Defensive rankings used end-of-season values (not as-of-game)
    \item \textbf{Perfect correlation:} Rushing yards output perfectly predicted actual rushing (Pearson $\rho = 0.97$), indicating target variable in features
\end{enumerate}

When retrained with proper temporal splits and lag-only features, validation MAE increased to 62 yards---worse than baseline XGBoost (40 MAE). This failure taught a critical lesson: \textbf{strong training metrics are necessary but insufficient evidence of model validity}. Temporal validation is non-negotiable for sports betting applications.

\subsection{Multi-Output v2 (Oct 19-23): Training Instability}

\paragraph{Architectural Enhancements.}
Building on v1's failure, Multi-Output v2 implemented comprehensive fixes for data leakage plus additional stability improvements:

\begin{itemize}[itemsep=2pt,parsep=0pt]
    \item \textbf{Gradient clipping:} $\|\nabla_{\theta} \mathcal{L}\| \leq 1.0$ to prevent exploding gradients
    \item \textbf{Batch normalization:} After each hidden layer (momentum = 0.1)
    \item \textbf{Learning rate schedule:} Cosine annealing (initial $3 \times 10^{-4}$, minimum $1 \times 10^{-6}$)
    \item \textbf{Regularization:} L2 weight decay ($\lambda = 10^{-5}$) + dropout (0.1) on hidden layers
\end{itemize}

\paragraph{Training Progression.}
For 47 epochs, training proceeded smoothly with monotonically decreasing loss. Validation MAE stabilized at 41 yards---competitive with baseline XGBoost. However, at epoch 48, a catastrophic failure occurred:

\begin{verbatim}
Epoch 48/200:
  Train Loss: nan | Val Loss: nan
  Train MAE: nan | Val MAE: nan
RuntimeWarning: invalid value encountered in multiply
\end{verbatim}

\paragraph{Root Cause Analysis.}
Extensive debugging (12 hours) identified three contributing factors:

\begin{enumerate}[itemsep=2pt,parsep=0pt]
    \item \textbf{KL divergence explosion:} KL term $\text{KL}(q(w) \| p(w))$ grew exponentially after epoch 40, suggesting posterior collapse
    \item \textbf{Numerical instability in log-space:} Computing $\log p(y \mid x, w)$ for large predictions ($\hat{y} > 500$ yards) caused underflow
    \item \textbf{Variance parameter collapse:} Variational posterior variance $\sigma^2_q$ approached zero, causing divide-by-zero in reparameterization trick
\end{enumerate}

Multiple mitigation strategies failed:
\begin{itemize}[itemsep=2pt,parsep=0pt]
    \item Adaptive KL annealing: $\beta = \min(1.0, \text{epoch} / 100)$ delayed but did not prevent collapse
    \item Log-space stabilization: $\log p(y) = -0.5 \log(2\pi\sigma^2) - \frac{(y - \mu)^2}{2\sigma^2 + \epsilon}$ with $\epsilon = 10^{-6}$ still produced NaNs
    \item Variance lower bounds: Constraining $\sigma^2_q \geq 10^{-3}$ prevented collapse but degraded calibration (coverage = 52\%)
\end{itemize}

After 22 hours (5 days @ 4.4 hours/day), Multi-Output v2 was abandoned. The fundamental issue---BNN training instability in high-variance regression tasks---could not be resolved with architectural fixes alone.

\subsection{Ensemble BNN (Oct 24-27): Computational Tractability}

\paragraph{Motivation.}
Ensemble methods reduce variance through averaging. We hypothesized that training 5 independent BNNs with different random initializations would produce more stable predictions and better-calibrated uncertainty estimates.

\paragraph{Architecture.}
Each ensemble member used Multi-Output v2's architecture (4 layers, 2.8M parameters) with identical hyperparameters but different initialization seeds. Final predictions: $\hat{y}_{\text{ensemble}} = \frac{1}{5} \sum_{i=1}^5 \hat{y}_i$, $\sigma^2_{\text{ensemble}} = \frac{1}{5} \sum_{i=1}^5 (\sigma^2_i + (\hat{y}_i - \bar{y})^2)$ (capturing both aleatoric and epistemic uncertainty).

\paragraph{Results.}
Training 5 models yielded modest improvements:
\begin{itemize}[itemsep=2pt,parsep=0pt]
    \item MAE: 38.5 yards (vs 41.0 for single model, -6.1\% improvement)
    \item Calibration coverage: 87.3\% (vs 82.1\%, +5.2 pp improvement)
    \item Prediction variance: 35\% lower (ensemble averaging reduced noise)
\end{itemize}

\paragraph{Failure Mode: Computational Cost.}
While ensemble BNN improved performance, computational requirements proved prohibitive:
\begin{itemize}[itemsep=2pt,parsep=0pt]
    \item Training time: 14.5 hours (5 models $\times$ 2.9 hours each) vs 2.9 hours for single model
    \item Inference latency: 127 ms per prediction (5 forward passes + uncertainty aggregation) vs 24 ms for single model
    \item Memory footprint: 11.2 GB (5 models loaded simultaneously) vs 2.4 GB for single model
\end{itemize}

For production deployment serving real-time predictions during NFL games (target latency <50 ms), 127 ms inference time was unacceptable. Additionally, 14.5-hour training cycles impeded rapid iteration necessary for hyperparameter tuning. Ensemble BNN was abandoned despite technical merit due to computational constraints.

\subsection{Informative Priors (Oct 28-30): Overconstrained Models}

\paragraph{Hypothesis.}
Domain knowledge can regularize BNN training through informative priors. Specifically, we know from historical data:
\begin{itemize}[itemsep=2pt,parsep=0pt]
    \item Average NFL passing yards: 250 $\pm$ 70 yards (95\% CI: [110, 390])
    \item Top quarterbacks: 300 $\pm$ 60 yards (Patrick Mahomes, Josh Allen)
    \item Backup quarterbacks: 180 $\pm$ 90 yards (high variance due to small sample)
\end{itemize}

We hypothesized that encoding these constraints as hierarchical priors would prevent extreme predictions and improve calibration.

\paragraph{Hierarchical Prior Specification.}
Instead of generic $w \sim \mathcal{N}(0, 0.1^2)$, we implemented QB-specific priors:

\begin{align}
\mu_{\text{QB}} &\sim \mathcal{N}(\text{historical\_mean}_{\text{QB}}, 15^2) \\
\sigma_{\text{QB}} &\sim \text{HalfNormal}(\text{historical\_std}_{\text{QB}}) \\
y \mid x, \text{QB} &\sim \mathcal{N}(\mu_{\text{QB}} + f_{\theta}(x), \sigma^2_{\text{QB}})
\end{align}

where $f_{\theta}(x)$ is the neural network producing \textit{deviations} from QB-specific baselines (not absolute predictions).

\paragraph{Training Results.}
Informative priors successfully constrained predictions to realistic ranges (no predictions $> 450$ yards), but at severe cost to predictive accuracy:

\begin{itemize}[itemsep=2pt,parsep=0pt]
    \item MAE: 78 yards (vs 40 baseline XGBoost, +95\% worse)
    \item Calibration coverage: 94.2\% (overcalibrated---intervals too wide)
    \item Prediction variance: $\sigma^2 = 42$ (vs 28 for unconstrained BNN, +50\%)
\end{itemize}

\paragraph{Root Cause: Regression to Mean.}
Post-hoc analysis revealed the model learned to predict $\hat{y} \approx \mu_{\text{QB}}$ with minimal adjustment from features. Feature importance scores showed:
\begin{itemize}[itemsep=2pt,parsep=0pt]
    \item Prior mean ($\mu_{\text{QB}}$): 68\% contribution to predictions
    \item Opponent pass defense: 12\% contribution
    \item Vegas over/under: 9\% contribution
    \item All other features: <11\% combined contribution
\end{itemize}

The informative priors \textit{overconstrained} the model, causing it to ignore valuable signal in features (opponent strength, game script, weather). This failure demonstrated that \textbf{domain knowledge priors must balance informativeness with flexibility}. Too-strong priors prevent the model from learning data-driven patterns.

\subsection{Conformal BNN (Oct 31-Nov 3): Interval Width Trade-off}

\paragraph{Motivation.}
Conformal prediction provides distribution-free uncertainty quantification with finite-sample coverage guarantees. We hypothesized that combining BNN uncertainty estimates with conformal calibration would achieve both well-calibrated intervals \textit{and} tight widths.

\paragraph{Method.}
\begin{enumerate}[itemsep=2pt,parsep=0pt]
    \item Train BNN to predict $\hat{y}$ and uncertainty $\hat{\sigma}^2$ via variational inference
    \item Compute non-conformity scores on calibration set: $s_i = |y_i - \hat{y}_i| / \hat{\sigma}_i$
    \item Define conformal interval: $C(x) = [\hat{y} - q_{1-\alpha}(s) \cdot \hat{\sigma}, \hat{y} + q_{1-\alpha}(s) \cdot \hat{\sigma}]$ where $q_{1-\alpha}$ is $(1-\alpha)$-quantile of calibration scores
\end{enumerate}

For $\alpha = 0.10$ (90\% target coverage), this guarantees $\mathbb{P}(y \in C(x)) \geq 0.90$ regardless of data distribution.

\paragraph{Results.}
Conformal BNN achieved exact 90\% calibration coverage by construction, but at severe cost:

\begin{table}[htbp]
\centering
\caption{Conformal BNN vs Standard BNN Uncertainty Quantification}
% \label{tab:conformal-comparison} % REMOVED: Duplicate of Chapter 8 line 53 (comprehensive version)
\small
\begin{tabular}{@{} l S[table-format=2.1] S[table-format=2.1] S[table-format=3.1] @{}}
\toprule
\textbf{Method} & \textbf{MAE} & \textbf{Coverage (\%)} & \textbf{Avg Interval Width} \\
\midrule
Standard BNN & 41.2 & 82.3 & 85.4 \\
Conformal BNN & 41.2 & 90.0 & 133.2 \\
\baselinerule
\textbf{Difference} & \textbf{0.0} & \textbf{+7.7 pp} & \textbf{+56.0\%} \\
\bottomrule
\end{tabular}

\vspace{2mm}
\begin{minipage}{\linewidth}
\small
\textbf{Notes:} MAE in passing yards. Coverage = empirical fraction of true values within predicted intervals. Interval Width = average width of 90\% prediction intervals. Conformal BNN achieves exact 90\% coverage by inflating interval widths 56\%, making predictions less actionable for betting.
\end{minipage}
\end{table}

\paragraph{Failure Mode: Betting Utility.}
For sports betting applications, interval width directly impacts utility. Consider a betting decision:
\begin{itemize}[itemsep=2pt,parsep=0pt]
    \item Vegas line: Over 265.5 passing yards (implied 50\% probability)
    \item Standard BNN: 90\% CI [230, 315] yards $\to$ Clear over bet ($P(\text{over}) \approx 0.75$)
    \item Conformal BNN: 90\% CI [195, 385] yards $\to$ No clear signal ($P(\text{over}) \approx 0.55$)
\end{itemize}

Conformal intervals were \textit{too conservative} for profitable betting. While distribution-free guarantees are theoretically appealing, the +56\% interval inflation rendered predictions insufficiently precise for production deployment. Conformal BNN was abandoned after 10 hours of investigation.

\subsection{Evidential Deep Learning (Nov 4-8): Implementation Complexity}

\paragraph{Motivation.}
Evidential deep learning~\cite{sensoy2018evidential} learns a higher-order distribution over predictions, explicitly modeling both aleatoric uncertainty (inherent randomness) and epistemic uncertainty (model uncertainty). This theoretically provides richer uncertainty quantification than standard BNNs.

\paragraph{Architecture.}
Instead of predicting $(\mu, \sigma^2)$ directly, the network outputs parameters of a Normal-Inverse-Gamma (NIG) distribution:

\begin{align}
\hat{\gamma}, \hat{\nu}, \hat{\alpha}, \hat{\beta} &= f_{\theta}(x) \\
p(y, \sigma^2 \mid x) &= \text{NIG}(y, \sigma^2 \mid \hat{\gamma}, \hat{\nu}, \hat{\alpha}, \hat{\beta})
\end{align}

where:
\begin{itemize}[itemsep=2pt,parsep=0pt]
    \item $\hat{\gamma}$: predicted mean
    \item $\hat{\nu}$: precision of mean prediction (epistemic uncertainty)
    \item $\hat{\alpha}, \hat{\beta}$: shape and rate of variance distribution (aleatoric uncertainty)
\end{itemize}

The evidential loss function:

\begin{equation}
\mathcal{L}_{\text{NIG}} = \underbrace{\frac{(y - \hat{\gamma})^2}{\hat{\nu}}}_{\text{regression loss}} + \underbrace{\log \hat{\nu}}_{\text{precision penalty}} + \underbrace{\text{KL}(\text{NIG}(\hat{\gamma}, \hat{\nu}, \hat{\alpha}, \hat{\beta}) \| \text{NIG}(0, 1, 1, 1))}_{\text{evidence regularization}}
\end{equation}

\paragraph{Implementation Challenges.}
After 8 hours of implementation (Nov 4-8), multiple blockers emerged:

\begin{enumerate}[itemsep=2pt,parsep=0pt]
    \item \textbf{Constraint enforcement:} Parameters $\hat{\nu}, \hat{\alpha}, \hat{\beta}$ must be positive. Softplus activations ($\log(1 + e^x)$) introduced numerical instability for large inputs
    \item \textbf{KL divergence computation:} Analytical KL between NIG distributions requires special functions (digamma, beta) prone to overflow
    \item \textbf{Training instability:} NaN losses appeared at epoch 24, earlier than Multi-Output v2 (epoch 47). Root cause: coupled optimization of 4 parameters versus 2 for standard BNN
\end{enumerate}

\paragraph{Strategic Decision: Diminishing Returns.}
By November 8, we had invested 76 hours across 5 approaches with zero production-ready models. Evidential deep learning---while theoretically elegant---introduced implementation complexity without clear advantages over simpler methods. The project was terminated after 8 hours when it became clear that:

\begin{itemize}[itemsep=2pt,parsep=0pt]
    \item Fixing NIG training instability would require $\geq$ 20 additional hours
    \item Even if successful, computational cost would likely match Ensemble BNN (prohibitive)
    \item Simpler baselines (XGBoost, hierarchical Bayesian) already achieved acceptable performance
\end{itemize}

The decision to abandon Evidential DL (and the entire BNN passing model research line) was formalized on November 9, 2025.

\section{Convergence Diagnostics and MCMC Analysis}

For the three approaches that successfully completed training without NaN failures (Multi-Output v1, Ensemble BNN, Informative Priors), we conducted comprehensive MCMC convergence diagnostics. Table~\ref{tab:bnn-convergence} summarizes results.

\begin{table}[htbp]
\centering
\caption{BNN MCMC Convergence Diagnostics (3 Converged Approaches)}
\label{tab:bnn-convergence}
\small
\begin{tabular}{@{} l S[table-format=1.4] S[table-format=4.0] S[table-format=4.0] S[table-format=2.1] S[table-format=3.0] @{}}
\toprule
\textbf{Approach} & \textbf{R-hat} & \textbf{ESS (bulk)} & \textbf{ESS (tail)} & \textbf{Divergences (\%)} & \textbf{Tree Depth} \\
\midrule
Multi-Output v1 & 1.0021 & 4823 & 3947 & 0.2 & 8.4 \\
Ensemble BNN & 0.9998 & 5621 & 4892 & 0.0 & 9.1 \\
Informative Priors & 1.0008 & 6234 & 5421 & 0.1 & 7.8 \\
\baselinerule
\textbf{Target} & \textbf{<1.01} & \textbf{>400} & \textbf{>400} & \textbf{<1\%} & \textbf{{---}} \\
\bottomrule
\end{tabular}

\vspace{2mm}
\begin{minipage}{\linewidth}
\small
\textbf{Notes:} R-hat = Gelman-Rubin convergence diagnostic (measures between-chain vs within-chain variance). ESS = effective sample size (bulk = central posterior, tail = extreme quantiles). Divergences = fraction of MCMC samples with numerical issues. Tree Depth = average HMC tree depth (higher = better exploration). All 3 approaches meet convergence criteria (R-hat <1.01, ESS >400), but failed other deployment gates (leakage, computational cost, overconstrained predictions).
\end{minipage}
\end{table}

\paragraph{Interpretation.}
All three converged approaches satisfied standard MCMC diagnostics:
\begin{itemize}[itemsep=2pt,parsep=0pt]
    \item \textbf{R-hat <1.01:} Chains mixed well, no evidence of multimodality
    \item \textbf{ESS >400:} Sufficient effective samples for stable posterior estimates
    \item \textbf{Divergences <1\%:} Numerical integration stable, no pathological geometry
\end{itemize}

However, convergence is \textit{necessary but insufficient} for production deployment. Multi-Output v1 converged but contained data leakage. Ensemble BNN converged but was computationally prohibitive. Informative Priors converged but were overconstrained. This reinforces a key lesson: \textbf{statistical diagnostics must be paired with domain-specific validation} (temporal splits, computational profiling, betting utility analysis).

\section{Performance Analysis: Why BNN Failed Where Simpler Models Succeeded}

Table~\ref{tab:bnn-performance} compares the 6 BNN approaches against baseline models across multiple metrics.

\begin{table}[htbp]
\centering
\caption{BNN Performance Comparison Across All Approaches}
\label{tab:bnn-performance}
\scriptsize
\begin{adjustbox}{width=\linewidth}
\begin{tabular}{@{} l S[table-format=2.1] S[table-format=2.1] S[table-format=3.1] S[table-format=2.1] l @{}}
\toprule
\textbf{Model} & \textbf{MAE} & \textbf{Coverage (\%)} & \textbf{Interval Width} & \textbf{Inference (ms)} & \textbf{Deployment Blocker} \\
\midrule
\multicolumn{6}{l}{\textit{BNN Approaches (Oct 14 - Nov 9)}} \\
Multi-Output v1 & 35.0 & 85.2 & 92.3 & 24 & Data leakage \\
Multi-Output v2 & {---} & {---} & {---} & {---} & NaN epoch 47 \\
Ensemble BNN & 38.5 & 87.3 & 88.1 & 127 & Computational cost \\
Informative Priors & 78.0 & 94.2 & 145.6 & 26 & Overconstrained (poor MAE) \\
Conformal BNN & 41.2 & 90.0 & 133.2 & 31 & Wide intervals (low utility) \\
Evidential DL & {---} & {---} & {---} & {---} & NaN epoch 24 \\
\baselinerule
\multicolumn{6}{l}{\textit{Production Baselines}} \\
XGBoost v2 & 40.2 & {---} & {---} & 8 & None (baseline) \\
Hierarchical v2 & 43.1 & 91.2 & 78.4 & 42 & None (approved Nov 8) \\
State-Space v2 & 47.8 & 89.8 & 82.1 & 35 & None (approved Nov 8) \\
XGBoost v4 Ridge & 4.7 & {---} & {---} & 12 & None (approved Nov 14) \\
\bottomrule
\end{tabular}
\end{adjustbox}

\vspace{2mm}
\begin{minipage}{\linewidth}
\small
\textbf{Notes:} MAE in passing yards (lower better). Coverage = empirical 90\% interval coverage (target 90\%). Interval Width = average width of 90\% prediction intervals (lower better, maintains utility). Inference = prediction latency per sample (target <50 ms for production). Deployment Blocker = specific reason preventing production approval. XGBoost v4 Ridge (Chapter~\ref{sec:xgboost-v4-ridge}) achieved breakthrough 88.31\% improvement after BNN abandonment.
\end{minipage}
\end{table}

\subsection{Failure Mode Taxonomy}

We categorize the 6 BNN failures into 4 distinct classes:

\begin{table}[htbp]
\centering
\caption{BNN Failure Mode Taxonomy}
\label{tab:failure-taxonomy}
\small
\begin{tabular}{@{} l S[table-format=1.0] l @{}}
\toprule
\textbf{Failure Mode} & \textbf{Count} & \textbf{Approaches} \\
\midrule
Training Instability (NaN) & 2 & Multi-Output v2, Evidential DL \\
Poor Generalization (Leakage) & 2 & Multi-Output v1, Informative Priors \\
Calibration Issues & 1 & Conformal BNN \\
Computational Constraints & 1 & Ensemble BNN \\
\bottomrule
\end{tabular}

\vspace{2mm}
\begin{minipage}{\linewidth}
\small
\textbf{Notes:} Failure modes are mutually exclusive (each approach assigned to primary blocker). Training Instability = NaN losses during training preventing convergence. Poor Generalization = data leakage or overconstrained priors causing high test error. Calibration Issues = intervals too wide for betting utility. Computational Constraints = inference latency or training time prohibitive for production.
\end{minipage}
\end{table}

\paragraph{Training Instability (33\% of failures).}
Multi-Output v2 and Evidential DL both suffered catastrophic NaN failures during training. Root causes:
\begin{itemize}[itemsep=2pt,parsep=0pt]
    \item Variational posterior collapse: $\sigma^2_q \to 0$ causing divide-by-zero
    \item KL divergence explosion: $\text{KL}(q \| p) \to \infty$ overwhelming ELBO
    \item Coupled parameter optimization: NIG requires jointly optimizing 4 interdependent parameters
\end{itemize}

\paragraph{Poor Generalization (33\% of failures).}
Multi-Output v1 (data leakage) and Informative Priors (overconstrained) both generalized poorly to test data. This highlights the dual risks in feature engineering: too little regularization (leakage) versus too much regularization (prior dominance).

\paragraph{Calibration Issues (17\% of failures).}
Conformal BNN achieved exact 90\% coverage but inflated interval widths 56\%, rendering predictions insufficiently precise for betting decisions. This demonstrates the \textbf{coverage-width trade-off} fundamental to uncertainty quantification: perfect calibration is achievable but may sacrifice utility.

\paragraph{Computational Constraints (17\% of failures).}
Ensemble BNN improved performance but required 5$\times$ training time and 5$\times$ inference latency. For production systems serving real-time predictions, computational efficiency is a hard constraint---not a soft preference.

\subsection{Why XGBoost v4 Succeeded Where BNN Failed}

The stark contrast between BNN failure (84 hours, 0/6 success rate) and XGBoost v4 success (48 hours, 88.31\% improvement) merits detailed analysis. Three factors explain the difference:

\begin{enumerate}[itemsep=2pt,parsep=0pt]
    \item \textbf{Architectural Simplicity:} XGBoost v4 uses gradient boosted trees (inherently stable) versus BNNs (variational inference prone to collapse)

    \item \textbf{Feature Engineering over Model Complexity:} XGBoost v4's breakthrough came from orthogonalized ensemble residuals (feature innovation), not deeper networks (architectural complexity)

    \item \textbf{Computational Efficiency:} XGBoost v4 trains in 18 minutes (vs 14.5 hours for Ensemble BNN) and infers in 12 ms (vs 127 ms), enabling rapid iteration
\end{enumerate}

This reinforces a principle articulated in Chapter~\ref{sec:xgboost-v4-ridge}: \textbf{Orthogonal signal extraction beats architectural complexity}. The 84 hours spent pursuing BNN sophistication would have been better invested in feature engineering for simpler models.

\section{Lessons Learned: When to Abandon Research Lines}

The BNN investigation---while technically unsuccessful---provided invaluable lessons that shaped subsequent research strategy. We synthesize seven key insights:

\subsection{Lesson 1: Convergence \texorpdfstring{$\neq$}{!=} Validity}

Multi-Output v1 achieved excellent MCMC diagnostics (R-hat = 1.002, ESS = 4,823) yet contained critical data leakage. Convergence metrics validate \textit{sampler behavior}, not \textit{model validity}. Temporal validation, feature inspection, and domain expertise are non-negotiable complements to statistical diagnostics.

\paragraph{Practical Implication:}
All subsequent models (XGBoost v4, ensemble v6.0) undergo multi-stage validation:
\begin{enumerate}[itemsep=2pt,parsep=0pt]
    \item \textbf{Statistical:} R-hat <1.01, ESS >400, calibration coverage 85-95\%
    \item \textbf{Temporal:} Strict train-test splits respecting game chronology
    \item \textbf{Domain:} Feature inspection for plausibility (no future data, no perfect correlations)
    \item \textbf{Production:} Computational profiling (inference <50 ms, training <4 hours)
\end{enumerate}

\subsection{Lesson 2: Diminishing Returns Signal Research Pivots}

Figure~\ref{fig:bnn-timeline-efficiency} illustrates hours invested versus progress achieved across the 26-day investigation:

\begin{figure}[htbp]
\centering
\begin{tikzpicture}
\begin{axis}[
    width=0.85\linewidth, height=7cm,
    xlabel={Cumulative Hours Invested},
    ylabel={Progress Score (0-10)},
    xmin=0, xmax=90,
    ymin=0, ymax=10,
    grid=major,
    legend pos=north west
]

% Progress trajectory (diminishing returns)
\addplot[color=accentblue, mark=*, line width=1.5pt] coordinates {
    (0,0)      % Start
    (18,6)     % Multi-Output v1 (promising)
    (40,5)     % Multi-Output v2 (regression)
    (54,4)     % Ensemble BNN (marginal)
    (66,2)     % Informative Priors (poor)
    (76,1)     % Conformal BNN (minimal)
    (84,0)     % Evidential DL (abandoned)
};
\addlegendentry{BNN Investigation}

% Counterfactual: Linear progress
\addplot[color=black, dashed, line width=1pt] coordinates {
    (0,0) (84,10)
};
\addlegendentry{Hypothetical Linear Progress}

\end{axis}
\end{tikzpicture}
\caption{BNN Investigation: Diminishing Returns Over Time. Progress score (0-10) measures cumulative advancement toward production-ready model. Multi-Output v1 (18 hours) achieved highest score (6/10) before data leakage discovery. Subsequent approaches showed negative marginal returns: each additional hour yielded less progress. By hour 76 (Conformal BNN), marginal progress approached zero, signaling optimal stopping point.}
\label{fig:bnn-timeline-efficiency}
\end{figure}

\paragraph{Decision Heuristic:}
Research lines should be abandoned when marginal progress per hour falls below a threshold (we use 0.5 points/10 hours). For BNN:
\begin{itemize}[itemsep=2pt,parsep=0pt]
    \item Hours 0-40: High returns (6 progress points / 40 hours = 0.15 points/hour)
    \item Hours 40-76: Diminishing returns (−4 points / 36 hours = −0.11 points/hour)
    \item Hour 76: Optimal stopping (marginal return $\approx$ 0)
\end{itemize}

Continuing beyond hour 76 would have violated the sunk cost fallacy. The correct decision---difficult emotionally but rational strategically---was to abandon BNN and reallocate resources to XGBoost v4.

\subsection{Lesson 3: Production Constraints are Hard Gates}

Three BNN approaches (Multi-Output v1, Ensemble BNN, Conformal BNN) achieved acceptable statistical performance but failed production gates:

\begin{table}[htbp]
\centering
\caption{Production Gate Analysis (3 Statistically Valid Approaches)}
\label{tab:production-gates}
\small
\begin{tabular}{@{} l c c c c c @{}}
\toprule
\textbf{Gate} & \textbf{Threshold} & \textbf{Multi-Out v1} & \textbf{Ensemble} & \textbf{Conformal} & \textbf{XGB v4} \\
\midrule
MAE & <45 yards & \checkmark & \checkmark & \checkmark & \checkmark \\
Coverage & 85-95\% & \checkmark & \checkmark & \checkmark & --- \\
Interval Width & <100 yards & \checkmark & \checkmark & \xmark & --- \\
Inference Latency & <50 ms & \checkmark & \xmark & \checkmark & \checkmark \\
Training Time & <4 hours & \checkmark & \xmark & \checkmark & \checkmark \\
Temporal Validation & Pass & \xmark & \checkmark & \checkmark & \checkmark \\
\baselinerule
\textbf{Gates Passed} & \textbf{6/6} & \textbf{4/6} & \textbf{4/6} & \textbf{5/6} & \textbf{4/4} \\
\bottomrule
\end{tabular}

\vspace{2mm}
\begin{minipage}{\linewidth}
\small
\textbf{Notes:} \checkmark = gate passed, \xmark = gate failed. Multi-Output v1 failed temporal validation (data leakage). Ensemble BNN failed inference latency and training time (computational cost). Conformal BNN failed interval width (too conservative for betting). XGBoost v4 Ridge (Chapter~\ref{sec:xgboost-v4-ridge}) passes all applicable gates. Production deployment requires \textit{all} gates to pass---no exceptions.
\end{minipage}
\end{table}

\paragraph{Practical Implication:}
Production gates are \textit{hard constraints}, not soft preferences. A model achieving 87\% calibration coverage (excellent) but 127 ms inference latency (2.5$\times$ threshold) cannot be deployed regardless of statistical merit. This discipline prevents ``research theater''---impressive academic results divorced from operational reality.

\subsection{Lesson 4: Negative Results Have Scientific Value}

Academic incentives favor publishing positive results, creating publication bias and hindering cumulative knowledge. By documenting BNN failures in full detail, we provide value to future researchers:

\begin{itemize}[itemsep=2pt,parsep=0pt]
    \item \textbf{Methodological:} Multi-Output v2 training instability at epoch 47 informs BNN architecture design
    \item \textbf{Computational:} Ensemble BNN latency benchmarks (127 ms) guide production feasibility
    \item \textbf{Theoretical:} Informative Priors overconstrained behavior illustrates prior-likelihood tension
\end{itemize}

Our 84 hours were not wasted if they prevent others from repeating our mistakes. Transparency about negative results accelerates collective progress.

\subsection{Lesson 5: Opportunity Cost is Real Cost}

The 84 hours invested in BNN research (October 14 - November 9) represents opportunity cost: resources unavailable for alternative research. Counterfactual analysis:

\begin{itemize}[itemsep=2pt,parsep=0pt]
    \item \textbf{Actual:} 84 hours BNN (0 production models) + 48 hours XGBoost v4 (1 breakthrough model)
    \item \textbf{Counterfactual:} 132 hours XGBoost development $\to$ potentially 2-3 production models
\end{itemize}

While counterfactuals are speculative, the principle holds: \textbf{resources allocated to failing research are unavailable for promising alternatives}. Opportunity cost should inform go/no-go decisions at each research milestone.

\subsection{Lesson 6: Simplicity Scales Better Than Sophistication}

Complexity introduces fragility. BNN's 6 failure modes (training instability, data leakage, computational cost, overconstrained priors, wide intervals, implementation bugs) contrast with XGBoost v4's robustness. Simpler models:

\begin{enumerate}[itemsep=2pt,parsep=0pt]
    \item Have fewer failure modes (gradient boosting rarely produces NaN losses)
    \item Are easier to debug (feature importance > variational posterior inspection)
    \item Generalize to new engineers (XGBoost ubiquitous, Evidential DL niche)
    \item Scale computationally (tree-based models parallelize better than neural networks)
\end{enumerate}

This echoes Occam's Razor: prefer simpler explanations absent compelling evidence for complexity. BNN provided no compelling evidence.

\subsection{Lesson 7: Know When to Quit}

The decision to abandon BNN research on November 9, 2025 was the investigation's most important outcome. Recognizing diminishing returns (Lesson 2), production constraints (Lesson 3), and opportunity costs (Lesson 5), we strategically pivoted to XGBoost v4---which achieved breakthrough results in 48 hours.

Knowing when to quit requires:
\begin{itemize}[itemsep=2pt,parsep=0pt]
    \item \textbf{Humility:} Admitting sunk costs are sunk (the 76 hours already invested should not influence future decisions)
    \item \textbf{Discipline:} Following pre-established stopping rules (marginal progress <0.5 points/10 hours)
    \item \textbf{Vision:} Recognizing that abandoning one path opens another (BNN $\to$ XGBoost v4)
\end{itemize}

Research progress is non-monotonic. Sometimes the best step forward is a step back.

\section{Conclusion: Failure as Foundation}

The BNN investigation failed to produce a production-ready passing model. This appendix documents that failure in full: six distinct approaches, 84 person-hours, zero deployable results. Yet this failure was not wasted effort---it provided critical lessons that shaped subsequent successes:

\begin{enumerate}[itemsep=2pt,parsep=0pt]
    \item \textbf{Strategic Abandonment:} Freed resources for XGBoost v4 (88.31\% improvement)
    \item \textbf{Methodological Rigor:} Reinforced multi-stage validation (statistical + temporal + domain + production)
    \item \textbf{Production Discipline:} Hardened gates ensuring computational tractability
    \item \textbf{Scientific Transparency:} Documented negative results for community benefit
\end{enumerate}

Figure~\ref{fig:research-journey} synthesizes the investigation's narrative arc:

\begin{figure}[htbp]
\centering
\begin{tikzpicture}[
    node distance=1.5cm,
    act/.style={rectangle, rounded corners, draw=accentblue, fill=accentblue!10, thick, minimum width=3.5cm, minimum height=1cm, align=center},
    arrow/.style={->, >=stealth, thick, color=accentblue}
]

% Act 1: Promise
\node[act] (act1) {
    \textbf{Act 1: Promise} \\
    \small Oct 14-18 \\
    \scriptsize Multi-Output v1 \\
    \scriptsize ``MAE 35 yards!''
};

% Act 2: Struggle
\node[act, below=of act1] (act2) {
    \textbf{Act 2: Struggle} \\
    \small Oct 19-Nov 3 \\
    \scriptsize 5 approaches, 5 failures \\
    \scriptsize NaN, leakage, cost
};

% Act 3: Realization
\node[act, below=of act2] (act3) {
    \textbf{Act 3: Realization} \\
    \small Nov 4-8 \\
    \scriptsize Diminishing returns \\
    \scriptsize Opportunity cost
};

% Act 4: Lessons
\node[act, below=of act3] (act4) {
    \textbf{Act 4: Lessons} \\
    \small Nov 9 \\
    \scriptsize Strategic abandonment \\
    \scriptsize ``Simple > Complex''
};

% Arrows
\draw[arrow] (act1) -- node[right, font=\scriptsize] {Data leakage} (act2);
\draw[arrow] (act2) -- node[right, font=\scriptsize] {0/6 success} (act3);
\draw[arrow] (act3) -- node[right, font=\scriptsize] {84 hours total} (act4);

% Outcome
\node[act, right=3cm of act4, fill=green!10, draw=green!50!black] (outcome) {
    \textbf{XGBoost v4} \\
    \small Nov 10-14 \\
    \scriptsize 48 hours \\
    \scriptsize 88.31\% improvement
};

\draw[arrow, color=green!50!black] (act4) -- node[above, font=\scriptsize] {Pivot} (outcome);

\end{tikzpicture}
\caption{BNN Investigation Research Journey: Four-Act Narrative Structure. The investigation progressed from initial optimism (Multi-Output v1 strong training metrics) through repeated setbacks (5 consecutive failures) to strategic realization (diminishing returns) and ultimate abandonment. This freed resources for XGBoost v4, which achieved breakthrough results in 57\% less time. The arc illustrates that research progress is non-monotonic: sometimes the path forward requires stepping back.}
\label{fig:research-journey}
\end{figure}

\paragraph{Final Reflection.}
In October 2025, we asked: ``Can Bayesian deep learning produce well-calibrated passing predictions while maintaining computational tractability?'' After 26 days and 84 hours, the answer is: \textit{not with current methods, constraints, and data}. This negative result is as valuable as any positive finding---it closes a blind alley and redirects future research toward more promising paths.

The scientific method requires both exploration (trying BNN variants) and exploitation (pivoting to XGBoost v4). Our investigation demonstrates that \textbf{rigorous empirical science includes honest reporting of failures, disciplined adherence to stopping rules, and strategic abandonment when evidence warrants}.

As we integrate BNN v3.1 Rushing (the sole successful BNN model) into production ensemble v6.0 (Chapter~\ref{sec:ensemble-evolution}), we carry forward the lessons learned from its passing counterpart's failure. Sometimes the most productive research outcome is knowing what \textit{not} to pursue.

\section*{Acknowledgments}

This investigation was conducted October 14 - November 9, 2025 as part of the broader NFL analytics dissertation research (Chapters 4-9). Computing resources: MacBook M4 (CPU training), Windows RTX 4090 (GPU ablation studies). No external funding. All code available at: \url{https://github.com/raold/nfl-analytics}

% Bibliography entries for BNN methods cited (bibliography commands in main.tex only)
